\honest{Human Evil and Muddled Thinking}{Eliezer Yudkowsky}{2007}

George Orwell saw country after country fall to totalitarianism, while ``the professors of free universities hailed the Soviet Union's purges as progress.'' You were born too late to remember it, because in your branch of time the Berlin Wall fell. If Orwell's name is not carved on one of its stones, it should be.

Orwell's weapon was clear writing. He knew ``that muddled language is muddled thinking.'' \nb{Orwell's own claim is that sloppy language ``makes it easier for us to have foolish thoughts.''} I quote him. Political speech is largely ``the defence of the indefensible'': British rule in India, the Russian purges, the atom bombs on Japan can be defended only by arguments too brutal to face, so politics speaks in euphemism. Bombing villages, machine-gunning cattle and burning huts is called ``PACIFICATION.'' If you simplify your English, you cannot speak the dialects of orthodoxy, and a stupid remark will be obvious ``even to yourself.'' That is the heart of \textsc{Overcoming Bias}.

We call the famines of Stalin and Mao evil because they were done ``by deliberate human intent.'' \nb{Whether they were intended is disputed. Standard accounts of the Chinese famine of 1959--1961 trace it to the Great Leap Forward, grain requisitioning and inflated harvest reports.} Evil must keep our revulsion latent, so human evil, ``wherever it exists,'' sets out to muddle thinking. \nb{Orwell's claim was about political language ``In our time.'' The post gives no evidence for the wider one.} In \textsc{1984} the villains falsify photographs, and O'Brien tortures Winston to make him say that two and two make five. I quote the scene.

I am ``continually aghast'' at intelligent people, such as Tyler Cowen, who do not think overcoming bias is important. Your mind separates you from an orangutan and built this world; its systematic failures matter. Would the Inquisition have tortured witches if everyone were an ideal Bayesian? \nb{Cowen's post questions ``a strong preference for overcoming bias, relative to other ends,'' calling bias ``one problem but not the only problem.'' It does not call it unimportant.} I quote Cowen: ``I view Robin's blog as exemplifying bias, and indeed showing that bias can be very useful.'' \nb{The sentence continues: ``especially if embedded in a broader discovery process with checks and balances.''} I hope this came only from trying to sound clever, and I ask whether Cowen really puts scope insensitivity on a level with plans to save lives. \nb{Cowen's post does not mention scope insensitivity.} Orwell fought a similar attitude in Stuart Chase, who came near to claiming that abstract words are meaningless and used this to excuse doing nothing: since you do not know what Fascism is, how can you fight it?

Perhaps overcoming bias does not look ``exciting enough'' without ``some clear evil to oppose.'' So I supply one: wherever there is evil, there are biases around it. ``The truth does have enemies.'' Had \textsc{Overcoming Bias} been a newsletter in the Soviet Union, all of us would have gone to labor camps.

Every great leap forward came from clearer thought, and every great woe, apart from a few natural catastrophes, from a stupidity. I give no evidence for either. ``Our last enemy is ourselves; and this is a war, and we are soldiers.'' \nb{Seven months earlier, in ``Politics Is the Mind-Killer,'' Yudkowsky had described how politics damages thinking: ``Arguments are soldiers.''}
